\chapter*{第 6 章 相似理论与量纲分析}\addcontentsline{toc}{chapter}{第 6 章 相似理论与量纲分析}\markboth{第 6 章 相似理论与量纲分析}{}

\begin{tishi}
本章是 818 的\textbf{小题富矿、大题边缘}章节：2014--2022 年真题中共出现约 \textbf{13 次}，
但\textbf{全部是选择、填空、判断、简答}（E1 slide33「第六章 相似理论与量纲分析＝量纲转化、相似理论」；
E2「第六章到第八章只有少量填空/判断/选择，掌握基本概念＋公式套用即可」）。\\
教材原书本章习题共 \textbf{12 题}（6.1--6.12，书 p116--118），本册\textbf{一题不漏}全部收录。\\
做题要求：概念题写全要点；比尺换算题必须写\textbf{已知 $\to$ 求 $\to$ 准则判据 $\to$ 比尺关系式 $\to$ 代入 $\to$ 结果 $\to$ 单位检验}；
量纲分析题必须写\textbf{量纲表 $\to$ 基本量纲个数 $m$ $\to$ $\pi$ 项个数 $n-m$ $\to$ 重复变量 $\to$ 指数方程 $\to$ 整理成准则}。\\
本章最重要的两条判据：\textbf{有自由液面 $\to$ 重力主导 $\to$ 弗劳德准则；封闭管道 $\to$ 黏滞力主导 $\to$ 雷诺准则}。
\end{tishi}

\begin{ti}
\sloppy
\emergencystretch=2em

\item 选择题。\xstarfull{5}\yiju{E4 2014 年试题选择第 9 题（讨论题，关于相似准则与量纲）；E4 2014 年试题选择第 10 题、2022 年单选第 9 题（"设计管道阀门实验首先要满足的相似准则是雷诺相似准则"，即本题第（1）小问）；E4 2016 年试题填空第 3、4 题（物理量无量纲组合，即本题第（5）（6）小问）；E4 2018 年试题选择第 5 题、2019 年试题选择第 5 题（$l$、$v$、$g$ 的无量纲组合 $v^{2}/(gl)$）；E4 2020 年试题选择第 7 题（相似准则的物理意义）；E1 slide33「第六章＝量纲转化、相似理论」}\nandu{中}
\begin{choices}
\item 设计有压管流模型实验时，应主要满足的相似准则是（\hspace{1cm}）。
\item[] A. 雷诺准则\quad B. 弗劳德准则\quad C. 欧拉准则\quad D. 柯西准则
\item 明渠水流模型实验应主要满足的相似准则是（\hspace{1cm}）。
\item[] A. 雷诺准则\quad B. 弗劳德准则\quad C. 欧拉准则\quad D. 韦伯准则
\item 压力输水管模型实验，长度比尺为 4，模型流量与原型流量之比应为（\hspace{1cm}）。
\item[] A. 2\quad B. 4\quad C. 8\quad D. 1/4
\item 明渠水流模型实验，长度比尺为 4，模型流量应为原型流量的（\hspace{1cm}）。
\item[] A. 1/2\quad B. 1/4\quad C. 1/8\quad D. 1/32
\item 物理量 $l$、$v$、$g$ 的无量纲组合是（\hspace{1cm}）。
\item[] A. $\dfrac{v}{g}$\quad B. $\dfrac{v^{2}}{gl}$\quad C. $\dfrac{gl}{v}$\quad D. $\dfrac{l}{v^{2}}$
\item 物理量 $p$（压强）、$\rho$（密度）、$l$（长度）、$Q$（流量）的无量纲组合是（\hspace{1cm}）。
\item[] A. $\dfrac{\rho Q^{2}}{pl}$\quad B. $\dfrac{p}{\rho Q}$\quad C. $\dfrac{p\rho l}{Q}$\quad D. $\dfrac{\rho Q}{pl}$\end{choices}

\item 两流动相似应满足的条件是什么？各条件的物理意义是什么？\xstarfull{5}\yiju{E4 2020 年试题判断第 8、9、10 题（相似三条件的"前提/核心/表现"表述、动力相似的主导作用）；E4 2018 年试题选择第 5 题（相似准则的物理意义）；E4 2014 年试题选择第 9 题；E1 slide33「第六章＝量纲转化、相似理论」；E2「第六章到第八章只有少量填空/判断/选择，掌握基本概念＋公式套用即可」}\nandu{易}

\item 写出下列物理量的量纲：速度、压强、密度、流量、力、加速度、运动黏度、动力黏度、力矩。\xstarfull{5}\yiju{E4 2014 年试题选择第 9 题（量纲式判断，含切应力量纲 $\mathrm{ML^{-1}T^{-2}}$）；E4 2020 年试题判断第 10 题（"量纲是物理量的实质，不受人为因素影响"）；E4 2019 年试题选择第 5 题；E1 slide33「第六章＝量纲转化」；教材式 6.37（E6）}\nandu{易}

\item 写出五个常用相似准数（弗劳德数、欧拉数、雷诺数、韦伯数、柯西数）的表达式，并说明各自的物理意义。\xstarfull{5}\yiju{E4 2019 年试题判断第 2 题（"弗劳德数的物理意义是惯性力和重力之比"）；E4 2020 年试题选择第 7 题、判断第 8、9、10 题（相似准则的物理意义与适用场合）；E4 2018 年试题选择第 5 题；E1 slide33「第六章＝相似理论」；教材式 6.42--6.46（E6）}\nandu{中}

\item 有一轿车，高 $h=1.5\ \mathrm{m}$，在公路上行驶，设计时速为 $108\ \mathrm{km/h}$。拟通过风洞中模型实验来确定此轿车在公路上以此速行驶时的空气阻力。已知该风洞实验段内气流的温度和轿车在公路上行驶时相同，轿车模型与原型的长度比尺 $k=\dfrac{2}{3}$，试求风洞实验段内的气流速度。\xstarfull{5}\yiju{E4 2022 年试题填空第 5 题（同题：轿车高 $1.5\ \mathrm{m}$、$108\ \mathrm{km/h}$、$k=2/3$，标准答案 $45\ \mathrm{m/s}$）；E4 2016 年试题填空第 4 题、2018 年试题选择第 5 题（模型相似换算）；E1 slide33「第六章＝相似理论」；教材式 6.35（E6）}\nandu{中}

\item 已知直径为 $15\ \mathrm{cm}$ 的输油管，油的流量为 $0.18\ \mathrm{m^{3}/s}$，油的运动黏度 $\nu_p=0.13\ \mathrm{cm^{2}/s}$。现用清水作模型实验，水的运动黏度 $\nu_m=0.013\ \mathrm{cm^{2}/s}$，为保证两流动相似，模型中水的流量只能取 $Q_m=0.018\ \mathrm{m^{3}/s}$。实验在 $5\ \mathrm{m}$ 长的模型管段上进行，试求：（1）模型的管径 $d_m$ 及长度比尺 $k_l$；（2）原型输油管 $100\ \mathrm{m}$ 长管段两端的压强差（用油柱高表示）。\xstarfull{3}\yiju{E4 2016 年试题第 7 题、2019 年试题第 3 题（同题改编：$d=15\ \mathrm{cm}$ 输油管的水模型实验）；教材例 6.1 为同一模型的例题方法（E6）；E2「第六章到第八章只有少量填空/判断/选择，掌握基本概念＋公式套用即可」，故取三星}\nandu{难}
\tuyi{习题 6.6 图（与教材例 6.1 同款装置）：上方为\textbf{原型输油管}，管径 $d_p=15\ \mathrm{cm}$，流量 $Q_p=0.18\ \mathrm{m^{3}/s}$，测段长 $100\ \mathrm{m}$；下方为\textbf{模型管}，用清水做实验流体，模型管长 $5\ \mathrm{m}$，流量 $Q_m=0.018\ \mathrm{m^{3}/s}$，两端接测压管（读得水柱差 $\Delta h_m$）。原型与模型之间按雷诺准则对应，压降按欧拉准则换算。}

\item 溢流坝泄流模型实验，原型最大泄流量 $Q_p=120\ \mathrm{m^{3}/s}$，实验室供水系统能提供的最大流量为 $Q_m=0.75\ \mathrm{m^{3}/s}$，试求此模型的最大长度比尺 $k_l$。若在模型上测得某处作用力 $F_m=2.8\ \mathrm{N}$，试求原型中相应位置的作用力 $F_p$。\xstarfull{3}\yiju{E4 2014 年试题第 8 题（溢流坝模型实验，参考答案 $F_p\approx1.13\ \mathrm{kN}$）；教材式 6.42 与弗劳德相似比尺（E6）；A6 指出"明渠流／堰流属旧大纲超纲内容"，2014 年后未再出，故取三星}\nandu{中}

\item 弧形闸门闸下泄流，长度比尺 $k_l=0.05$，模型实验测出下游收缩断面平均流速为 $2\ \mathrm{m/s}$，流量为 $35\ \mathrm{L/s}$，闸门上的总压力为 $40\ \mathrm{N}$，试求原型中相应的下游收缩断面平均流速、流量和闸门上的总压力。\xstarfull{4}\yiju{E4 2021 年试题第 6 题（闸下泄流模型实验，弗劳德相似换算流速、流量、总压力）；E4 2014 年试题第 8 题（同类比尺换算）；教材式 6.42--6.47（E6）}\nandu{中}

\item 恒定有压管流的下临界流速 $u_c$ 与管径 $d$、流体的动力黏度 $\mu$ 和管内流体的密度 $\rho$ 有关，试用瑞利法求临界雷诺数 $\Rey_c$ 的表达式。\xstarfull{5}\yiju{E4 2019 年试题判断第 8 题（"用瑞利法推求物理过程的方程式，待求的量纲指数不能超过 4 个"，考的就是瑞利法的适用范围，与本题同一考点）；E4 2014 年试题选择第 9、10 题（量纲转化）；E1 slide33「第六章＝量纲转化」；教材 6.2.3 节与教材式 6.52--6.55（E6）}\nandu{中}

\item 已知水泵的轴功率 $N$ 与泵轴的转矩 $M$、角速度 $\omega$ 有关，试用瑞利法求轴功率 $N$ 的表达式。\xstarfull{4}\yiju{E4 2019 年试题判断第 8 题（瑞利法的适用范围）；E4 2014 年试题选择第 9 题（量纲式判断）；E1 slide33「第六章＝量纲转化」；教材 6.2.3 节（E6）}\nandu{易}

\item 影响液体边壁切应力 $\tau_w$ 的因素有断面平均流速 $v$、水力半径 $R$、管壁粗糙度 $\Delta$、液体密度 $\rho$ 和动力黏度 $\mu$，试用 $\pi$ 定理导出
\[ \frac{\tau_w}{\rho v^{2}}=f\!\left(\Rey,\ \frac{\Delta}{R}\right) \]
\xstarfull{5}\yiju{E4 2014 年试题选择第 9、10 题（量纲与相似准则）；E4 2019 年试题选择第 5 题、2020 年试题选择第 7 题（$\pi$ 定理与无量纲组合的命题方式）；E1 slide33「第六章＝量纲转化」；教材 6.2.3 节与教材式 6.56--6.60（E6）}\nandu{难}

\item 水流围绕桥墩流动，绕流阻力 $F$ 与桥墩宽度 $b$、水流速度 $v$、水的密度 $\rho$、动力黏度 $\mu$ 和重力加速度 $g$ 有关，试用 $\pi$ 定理导出
\[ F=\rho b^{2}v^{2}f\!\left(\Rey,\ \Fr\right) \]
\xstarfull{5}\yiju{E4 2019 年试题选择第 5 题、2020 年试题选择第 7 题（绕流阻力与 $\pi$ 定理的无量纲组合命题方式）；E4 2014 年试题选择第 9、10 题；E1 slide33「第六章＝量纲转化、相似理论」；教材 6.2.3 节与教材式 6.62--6.74（E6）}\nandu{难}

\end{ti}

\begin{tishi}
\textbf{本章 12 题的真题命中情况（按 2014--2022 年 818 试卷统计）：}\\
{\footnotesize
\begin{tabularx}{\textwidth}{@{}llX@{}}
\toprule
题号 & 教材页 & 对应真题（原样或改编） \\
\midrule
6.1 & 书 p116 & 2014 选 9/10；2016 填 3/4；2018 选 5；2019 选 5、判 2；2020 选 7、判 8/9/10；2021 简答 3；2022 单选 9 \\
6.2 & 书 p117 & 2020 判 8/9/10；2014 选 9（概念） \\
6.3 & 书 p117 & 2014 选 9；2020 判 10；2019 选 5（量纲式） \\
6.4 & 书 p117 & 2019 判 2；2020 选 7、判 9；2018 选 5 \\
6.5 & 书 p117 & 2022 填 5（标准答案 $45\ \mathrm{m/s}$） \\
6.6 & 书 p117 & 2016 第 7 题；2019 第 3 题 \\
6.7 & 书 p117 & 2014 第 8 题（$F_p\approx1.13\ \mathrm{kN}$） \\
6.8 & 书 p118 & 2021 第 6 题 \\
6.9 & 书 p118 & 2019 判 8（瑞利法适用范围）；2014 选 9/10 \\
6.10 & 书 p118 & 2019 判 8；2014 选 9 \\
6.11 & 书 p118 & 2014 选 9/10；2019 选 5；2020 选 7 \\
6.12 & 书 p118 & 2019 选 5；2020 选 7；2014 选 9/10 \\
\bottomrule
\end{tabularx}}
\textbf{复习优先级：}6.1、6.4 的准则判据（必背）$>$ 6.3 的量纲表与 6.9--6.12 的量纲分析六步（必会）
$>$ 6.5--6.8 的比尺换算（会算）$>$ 6.2 的概念表述（速答）。
\end{tishi}
